%%%%%%%%%%%%%%%%%%%%%%%%%%%%%%%%%
%
%       EXERCÍCIO
%
%%%%%%%%%%%%%%%%%%%%%%%%%%%%%%%%%

\definevalorquestao

\begin{exercicioBanco}[\valorquestao]
\noindent
Em um processo de controle de qualidade na indústria química, analisam-se o número de impurezas metálicas ($X$) e o número de defeitos de superfície ($Y$) por lote produzido. A distribuição de probabilidade conjunta de $(X,Y)$ é dada por:
\[
\renewcommand{\arraystretch}{1.5}
\begin{array}{|c|c|c|c|}
\hline
X \backslash Y & 0 & 1 & 2 \\ \hline
0 & 1/12 & 1/6 & 1/12 \\ \hline
1 & 1/6 & 1/4 & 1/12 \\ \hline
2 & 1/12 & 1/12 & 0 \\ \hline
\end{array}
\]

\begin{enumerate}[(a)]
    \item Calcule $P(X \ge 1, Y \le 1)$ e $P(X = Y)$.
    \item Determine $E(X)$, $E(Y)$ e $\text{Cov}(X,Y)$.
    \item As variáveis $X$ e $Y$ são independentes? Justifique.
\end{enumerate}
\end{exercicioBanco}

